\section{Qualitative Case Study on Hallucination Reduction}\label{app:case_study}
We present a qualitative case study on object hallucination using the MSCOCO dataset~\citep{Chen2015MicrosoftCC}. Captions generated using LLaVA-OV~\citep{li2024llava} and using our proposed \algname are compared. Text in red indicates objects that are not present in the original image but are generated by the model. According to the case study, LLaVA-OV~\citep{li2024llava} generated hallucinated objects, \textit{e.g.}, 'sink', 'bird'. Our proposed \algname demonstrates superior performance by generating fewer object hallucinations. Additionally, it produces longer and more detailed captions compared to LLaVA-OV~\citep{li2024llava}.

\begin{tcolorbox}[colback=blue!5!white, colframe=blue!50!black, width=\textwidth, sharp corners, boxrule=0.5mm]
    % \begin{tabular}{p{0.3\textwidth} p{0.6\textwidth}} % Define two columns with fixed widths
        \centering % Center the image in its column
        \includegraphics[width=0.35\linewidth]{fig/000000191013.jpg} 
        % &
        \vspace{0.5em}
    \noindent
    \begin{tikzpicture}
        \draw[dashed] (0,0) -- (\textwidth,0); % 定义虚线
    \end{tikzpicture}
    \vspace{0.5em}
        \parbox[c]{\linewidth}{ % Center-align text vertically and horizontally
            \textbf{\textit{LLaVA-OV~\citep{li2024llava} Generated Caption:}} In the heart of a bustling city, a large, vibrant mural adorns the side of a brick building. The mural is a lively depiction of a yellow \textcolor{green}{clock} with black hands and numbers, standing out against the red background. The \textcolor{green}{clock} is not alone; it's accompanied by a red \textcolor{red}{bird} perched on its top and a gold trumpet resting on its bottom, adding an element of whimsy to the scene.
            The building itself is situated on a busy street corner. A yellow taxi cab and a red van are parked in front of it, hinting at the hustle and bustle of city life. Above, the sky is a clear blue, dotted with white clouds, providing a serene backdrop to this urban landscape.
            The image captures the essence of city life - the blend of art and architecture, the constant movement of vehicles, and the ever-changing sky. It's a snapshot of a moment in time in a city that's always on the move.
            \\
            \textcolor{blue}{\textbf{\algname (Ours) Generated caption:}} The image captures a large, ornate \textcolor{green}{clock} mounted on the side of a building. The \textcolor{green}{clock} face is circular and prominently features Roman numerals for the hours, with the hands indicating a time of approximately 12:30. Above the \textcolor{green}{clock} face, there is a decorative element that includes a red shield with the Vodafone logo in white, flanked by two golden trumpets on either side. Below the \textcolor{green}{clock} face, there are two banners; the upper banner reads "İşte ben de bir kral olur" which translates to "I too want to be a king" in Turkish, while the lower banner states "Vodafone 24 saat hizmet" meaning "Vodafone 24-hour service."
            In the foreground, there are several vehicles parked along the street, including a yellow taxi cab, a white van, and a maroon SUV. The sky above is clear with a few clouds, suggesting it might be late afternoon or early evening based on the lighting.
            The building itself appears to be constructed with red brick, and there is another advertisement visible on the right side of the building, though it is partially obscured by the \textcolor{green}{clock} structure. The overall style of the image is a real-life photograph taken during daylight hours.
        }
    % \end{tabular}
\end{tcolorbox}

\begin{tcolorbox}[colback=blue!5!white, colframe=blue!50!black, width=\textwidth, sharp corners, boxrule=0.5mm]
    % \begin{tabular}{p{0.3\textwidth} p{0.6\textwidth}} % Define two columns with fixed widths
        \centering % Center the image in its column
        \includegraphics[width=0.4\linewidth]{fig/000000237928.jpg} 
        % &
        \vspace{0.5em}
    \noindent
    \begin{tikzpicture}
        \draw[dashed] (0,0) -- (\textwidth,0); % 定义虚线
    \end{tikzpicture}
    \vspace{0.5em}
        \parbox[c]{\linewidth}{ % Center-align text vertically and horizontally
            \textbf{\textit{LLaVA-OV~\citep{li2024llava} Generated Caption:}} This image captures a modern kitchen bathed in soft light. The dominant color scheme is white, punctuated by the sleek silver of the appliances and the vibrant red of the kettle. 
            The kitchen is equipped with a stainless steel \textcolor{green}{refrigerator} standing tall on the right, its surface reflecting the ambient light. Adjacent to it, a black stove waits ready for culinary adventures, while a \textcolor{green}{microwave} sits atop it, offering convenience at arm's reach.
            The heart of the kitchen is a \textcolor{red}{sink} nestled between the stove and a window. The window allows natural light to filter in, illuminating the space and creating a warm atmosphere. Above the sink, white cabinets offer ample storage space, their pristine surfaces adding to the overall clean aesthetic of the kitchen.
            A red kettle rests on the counter next to the sink, its bold color contrasting beautifully with the surrounding whites. Its presence suggests many mornings started with a comforting cup of tea.
            Overall, this kitchen exudes a sense of modern elegance and functionality, with every element thoughtfully placed for ease of use and visual appeal.
            \\
            \textcolor{blue}{\textbf{\algname (Ours) Generated caption:}} The image depicts a modern kitchen with a clean and organized appearance. The kitchen features white cabinetry that includes both upper and lower cabinets, providing ample storage space. The upper cabinets have two doors each, while the lower cabinets have one door each, all fitted with silver handles.
            On the countertop, there is a stainless steel \textcolor{green}{microwave} \textcolor{green}{oven} positioned to the left of a red kettle. The \textcolor{green}{microwave} has a digital display and control panel on its front side. Next to the \textcolor{green}{microwave}, there is a red kettle with a black handle, which adds a pop of color to the otherwise neutral-toned kitchen.
            To the right of the \textcolor{green}{microwave} and kettle, there is a stainless steel \textcolor{green}{refrigerator} with a freezer compartment at the bottom. The \textcolor{green}{refrigerator} has a sleek design with a handle located on the right side for accessing the freezer and a separate handle for the main compartment.
            Below the countertop, there is a built-in \textcolor{green}{oven} with a glass window and a red \textcolor{green}{oven} mitt hanging from the \textcolor{green}{oven} handle. The \textcolor{green}{oven} appears to be part of the cabinetry, seamlessly integrated into the design.
            The backsplash behind the countertop consists of dark red tiles, which contrast with the white cabinetry and add depth to the space. The floor is covered with light-colored carpeting, which complements the overall color scheme of the kitchen.
            The lighting in the kitchen is soft and ambient, coming from recessed lights in the ceiling, which creates a warm and inviting atmosphere. There are no visible windows or other sources of natural light in the image.
            Overall, the kitchen is designed with functionality and aesthetics in mind, featuring modern appliances and a cohesive color palette that combines white, silver, and red accents.
            \\  
        }
    % \end{tabular}
    
    % \vspace{0.5em}
    % \noindent
    % \begin{tikzpicture}
    %     \draw[dashed] (0,0) -- (\textwidth,0); % 定义虚线
    % \end{tikzpicture}
    % \vspace{0.5em}
    
\end{tcolorbox}

% \section{Zero-shot and Linear Probe Classification}

% We present additional results on zero-shot and linear probe classification over 7 datasets in \tabref{tab:classification}. For all the experiments, we follow the setup of CLIP~\citep{radford2021learning}. \algname consistently demonstrates comparable performance to CLIP.

% % \setlength{\tabcolsep}{2pt}
% \begin{table}[h]
%     \centering
%     \setlength{\tabcolsep}{2pt} 
%     \renewcommand{\arraystretch}{1} 
%     \resizebox{\linewidth}{!}{ 
%     \begin{tabular}{lcccccccccccccc}
%         \toprule
%         & \multicolumn{7}{c}{Zero-Shot} & \multicolumn{7}{c}{Linear Probe} \\
%         \midrule
%         Model & CIFAR10 & CIFAR100 & MNIST & STL10 & SUN397 & SST2 & GTSRB &  CIFAR10 & CIFAR100 & MNIST & STL10 & SUN397 & SST2 & GTSRB \\
%         \midrule
%         CLIP~\citep{radford2021learning} (ViT-B/32) & 89.8 & \textbf{65.1} & 48.2 & \textbf{97.1} &  \textbf{63.2} & 59.6 & 32.2 & 95.1 & 80.5 & \textbf{99.0} & 98.3 & 76.6 & 70.8 & 85.3 \\
%         \algname-0.5B & 90.5 & 60.0 & \textbf{48.4} & 93.9  & 38.0 & \textbf{60.3} & \textbf{45.6} & 97.0 & \textbf{82.7} & 98.1 & 99.0 & 76.0 & 76.3 & \textbf{87.6} \\
%         \algname-7B & \textbf{93.1} & 63.7  & 33.0 & 96.9 & 56.7 & 54.1 & 43.6 & \textbf{97.4} & 82.0 & 98.6 & \textbf{99.1} & \textbf{76.9} & \textbf{86.6} & 87.4 \\
%         \bottomrule
%     \end{tabular}
%     }
%     \caption{Zero-shot and linear probe classification results. \algname achieves comparable performance to CLIP~\citep{radford2021learning}. }
%     \label{tab:classification}
% \end{table}